\section{Boundary examples and non-implications}\label{sec:boundaries}
The following examples specify the distinctions used in the theorems. They
are part of the mathematical scope, not objections to studying stronger
results under stronger hypotheses.

\begin{example}[Less statistical information, more exact prediction labels]
\label{ex:bit-gaussian}
Let $Z\sim\operatorname{Bernoulli}(1/2)$ and initially observe $X=Z$.
The future target is again $Z$. Two labels retain the posterior mean exactly.
Garble $X$ into $Y=X+\xi$ with independent $\xi\sim N(0,1)$. Bayes' formula
and the ratio of the two Gaussian densities give
\[
 p(Y)=\Pp(Z=1\mid Y)=\frac1{1+e^{1/2-Y}}.
\]
This is a smooth strictly increasing bijection from $\R$ to $(0,1)$.
The density of $Y$ is positive everywhere, so the posterior law is
continuous and has a positive density on every compact subinterval of
$(0,1)$. No finite set of centres can equal this posterior almost surely.
More quantitatively, restrict to a compact interval where the density has
a positive lower bound and apply Lemma~\ref{lem:attained-chart} in one
dimension. Every fixed $M$ has strictly positive optimal excess squared
error. Both oracles are different: absolute decision risk obeys statistical
data processing, but excess coding risk relative to each oracle need not.
\end{example}

\begin{example}[The same mean map with two noise geometries]
Let $Z\sim N(\lambda,1)$. For known $\epsilon>0$, observing
$Y_\epsilon=\epsilon Z$ is equivalent to observing $Z$ by division, and
its Fisher information for $\lambda$ is one. Observing
$\widetilde Y_\epsilon=\epsilon\lambda+\xi$ with independent
$\xi\sim N(0,1)$ has Fisher information $\epsilon^2$. Both means equal
$\epsilon\lambda$. Only the first covariance is scaled with the mean.
Consequently a small singular value of a mean map does not determine its
statistical information without the noise law.
\end{example}

\begin{example}[Rank does not record global fibres or contact order]
The derivative of $x\mapsto x^2$ is nonzero at $1$ and $-1$, yet their
images coincide. For another distinction, let $Y_i=\lambda^r+\xi_i$,
$\lambda\ge0$, with independent unit Gaussian noise and $r\ge2$. Every
mean map has zero derivative at zero, but its local amplitude scale is
$n^{-1/(2r)}$. Indeed the divergence between $\lambda=0$ and a positive
alternative is $n\lambda^{2r}/2$. Taking this bounded gives the lower
resolution scale by two-point testing. For an upper bound use
$\widehat\lambda=(\overline Y_+)^{1/r}$ and
$|x_+^{1/r}-b^{1/r}|\le|x-b|^{1/r}$ for $b\ge0$; its mean squared
error is $O(n^{-1/r})$. The order of contact, not just the rank, matters.
\end{example}

\begin{example}[Reachability is not attained mass]
A protocol may permit all commands in an interval while the exploration
law always chooses a single command. The set of legally attainable
predictions can be continuous even if that particular exploration law
visits only one prediction. Another exploration on the same command
interface can spread positive mass over a full interval. Their average
quantization risks differ. Theorem~\ref{thm:refresh} avoids this ambiguity
by specifying independent uniform commands and deriving their density.
\end{example}

\begin{example}[Equal first two moments do not give equal laws]
The two laws
$\frac12\delta_{-1}+\frac12\delta_1$ and
$\frac12\delta_0+\frac14\delta_{-\sqrt2}+\frac14\delta_{\sqrt2}$
have mean zero and variance one, but are different. Equality of a finite
jet or covariance needs a model-specific identifiability theorem to imply
anything stronger.
\end{example}

\begin{example}[Fisher preservation need not be sufficiency]
The Kagan--Shepp phenomenon, analyzed in \cite{Pollard2013}, gives an
insufficient statistic preserving Fisher information. We do not reproduce
that external construction here or claim it as new. Its precise logical
role is to prevent promotion of the equality case of
Theorem~\ref{thm:score} into the parameter-uniform recovery assertion of
Theorem~\ref{thm:entropy-gap}.
\end{example}

\begin{example}[Small variation error can change the exponential rate]
For $a_n\to\infty$ and $c>0$, let
\[
 P_n=\delta_0,\qquad
 Q_n=(1-e^{-a_nc})\delta_0+e^{-a_nc}\delta_1.
\]
Their total variation is $e^{-a_nc}\to0$. At speed $a_n$, however,
$P_n$ gives rate $+\infty$ to $1$, while $Q_n$ gives it rate $c$.
Thus even exponentially small variation error need not preserve a complete
rate function. Appropriate superexponential approximation hypotheses are
substantially stronger than an ordinary vanishing simulation defect.
\end{example}

\begin{example}[Linear-source pressure can convexify a nonconvex rate]
Let $P_n$ put probability $1/2$ at each of $-1$ and $1$, for all $n$.
At speed $n$ its rate vanishes at those two points and is infinite elsewhere.
Its limiting pressure for linear source $sx$ is $|s|$. The Legendre
transform of $|s|$ vanishes on the entire interval $[-1,1]$.
The full bounded-continuous-source Laplace functional and the restricted
linear-source transform are not interchangeable.
\end{example}

\begin{example}[Observation does not ensure Markov closure]
Consider the deterministic chain $a\mapsto a$, $b\mapsto c$, $c\mapsto c$,
with report $q(a)=q(b)=0$, $q(c)=1$. The current report $0$ does not specify
the next report: the two compatible hidden states give different answers.
Thus no single reduced kernel works for every preparation. Future-test
equivalence distinguishes $a$ and $b$; merely keeping their current report
does not.
\end{example}

\begin{example}[Zero evidence has no canonical normalized limit]
On a two-point space the positive measures $(\delta,0)$ and $(0,\delta)$
both tend to the zero measure as $\delta\downarrow0$. Their normalized
probabilities tend to different point masses. A continuous normalized
extension at zero cannot be assigned from the zero measure alone.
\end{example}

\begin{example}[A simulator may need delay memory]
Let the unknown parameter be $\lambda\in\{0,1\}$. A source with no
control actions reveals $\lambda$ at time one and constants thereafter.
A target reveals a constant at time one and $\lambda$ at time two.
Their terminal information is equivalent. However, a parameter-independent
simulator with only one retained state after time one, and no other retained
tape, has the same internal state and the same new inputs under both
parameters. Its second output therefore has the same law under both
parameters and cannot equal the required two different point masses.
Two states suffice by retaining the bit. If a target controller also
retains its own state, the direct implementation is a product. Terminal
sufficiency does not erase this state cost.
\end{example}

\begin{example}[A finite command alphabet is a different experiment]
At a fixed finite time, finite actions and finite reports produce only
finitely many histories. With enough labels they can be retained exactly.
The positive lower bound of Theorem~\ref{thm:refresh} for all finite $M$
uses the stipulated continuously distributed commands. Rounding the
physical command alphabet itself changes that experiment. Approximating
only the filter's arithmetic while observing the original experiment is
instead the perturbation problem in Section~\ref{sec:emulation}.
\end{example}

These examples rule out four unqualified shortcuts: deriving global
identifiability from local rank, deriving memory from information order,
deriving statistical rates from a noiseless feature map, and deriving
exponential asymptotics from ordinary variation approximation. The positive
results in this paper specify the extra structure needed in each case.
